%% ODER: format ==         = "\mathrel{==}"
%% ODER: format /=         = "\neq "
%
%
\makeatletter
\@ifundefined{lhs2tex.lhs2tex.sty.read}%
  {\@namedef{lhs2tex.lhs2tex.sty.read}{}%
   \newcommand\SkipToFmtEnd{}%
   \newcommand\EndFmtInput{}%
   \long\def\SkipToFmtEnd#1\EndFmtInput{}%
  }\SkipToFmtEnd

\newcommand\ReadOnlyOnce[1]{\@ifundefined{#1}{\@namedef{#1}{}}\SkipToFmtEnd}
\usepackage{amstext}
\usepackage{amssymb}
\usepackage{stmaryrd}
\DeclareFontFamily{OT1}{cmtex}{}
\DeclareFontShape{OT1}{cmtex}{m}{n}
  {<5><6><7><8>cmtex8
   <9>cmtex9
   <10><10.95><12><14.4><17.28><20.74><24.88>cmtex10}{}
\DeclareFontShape{OT1}{cmtex}{m}{it}
  {<-> ssub * cmtt/m/it}{}
\newcommand{\texfamily}{\fontfamily{cmtex}\selectfont}
\DeclareFontShape{OT1}{cmtt}{bx}{n}
  {<5><6><7><8>cmtt8
   <9>cmbtt9
   <10><10.95><12><14.4><17.28><20.74><24.88>cmbtt10}{}
\DeclareFontShape{OT1}{cmtex}{bx}{n}
  {<-> ssub * cmtt/bx/n}{}
\newcommand{\tex}[1]{\text{\texfamily#1}}	% NEU

\newcommand{\Sp}{\hskip.33334em\relax}


\newcommand{\Conid}[1]{\mathit{#1}}
\newcommand{\Varid}[1]{\mathit{#1}}
\newcommand{\anonymous}{\kern0.06em \vbox{\hrule\@width.5em}}
\newcommand{\plus}{\mathbin{+\!\!\!+}}
\newcommand{\bind}{\mathbin{>\!\!\!>\mkern-6.7mu=}}
\newcommand{\rbind}{\mathbin{=\mkern-6.7mu<\!\!\!<}}% suggested by Neil Mitchell
\newcommand{\sequ}{\mathbin{>\!\!\!>}}
\renewcommand{\leq}{\leqslant}
\renewcommand{\geq}{\geqslant}
\usepackage{polytable}

%mathindent has to be defined
\@ifundefined{mathindent}%
  {\newdimen\mathindent\mathindent\leftmargini}%
  {}%

\def\resethooks{%
  \global\let\SaveRestoreHook\empty
  \global\let\ColumnHook\empty}
\newcommand*{\savecolumns}[1][default]%
  {\g@addto@macro\SaveRestoreHook{\savecolumns[#1]}}
\newcommand*{\restorecolumns}[1][default]%
  {\g@addto@macro\SaveRestoreHook{\restorecolumns[#1]}}
\newcommand*{\aligncolumn}[2]%
  {\g@addto@macro\ColumnHook{\column{#1}{#2}}}

\resethooks

\newcommand{\onelinecommentchars}{\quad-{}- }
\newcommand{\commentbeginchars}{\enskip\{-}
\newcommand{\commentendchars}{-\}\enskip}

\newcommand{\visiblecomments}{%
  \let\onelinecomment=\onelinecommentchars
  \let\commentbegin=\commentbeginchars
  \let\commentend=\commentendchars}

\newcommand{\invisiblecomments}{%
  \let\onelinecomment=\empty
  \let\commentbegin=\empty
  \let\commentend=\empty}

\visiblecomments

\newlength{\blanklineskip}
\setlength{\blanklineskip}{0.66084ex}

\newcommand{\hsindent}[1]{\quad}% default is fixed indentation
\let\hspre\empty
\let\hspost\empty
\newcommand{\NB}{\textbf{NB}}
\newcommand{\Todo}[1]{$\langle$\textbf{To do:}~#1$\rangle$}

\EndFmtInput
\makeatother
%
%
%
%
%
%
% This package provides two environments suitable to take the place
% of hscode, called "plainhscode" and "arrayhscode". 
%
% The plain environment surrounds each code block by vertical space,
% and it uses \abovedisplayskip and \belowdisplayskip to get spacing
% similar to formulas. Note that if these dimensions are changed,
% the spacing around displayed math formulas changes as well.
% All code is indented using \leftskip.
%
% Changed 19.08.2004 to reflect changes in colorcode. Should work with
% CodeGroup.sty.
%
\ReadOnlyOnce{polycode.fmt}%
\makeatletter

\newcommand{\hsnewpar}[1]%
  {{\parskip=0pt\parindent=0pt\par\vskip #1\noindent}}

% can be used, for instance, to redefine the code size, by setting the
% command to \small or something alike
\newcommand{\hscodestyle}{}

% The command \sethscode can be used to switch the code formatting
% behaviour by mapping the hscode environment in the subst directive
% to a new LaTeX environment.

\newcommand{\sethscode}[1]%
  {\expandafter\let\expandafter\hscode\csname #1\endcsname
   \expandafter\let\expandafter\endhscode\csname end#1\endcsname}

% "compatibility" mode restores the non-polycode.fmt layout.

\newenvironment{compathscode}%
  {\par\noindent
   \advance\leftskip\mathindent
   \hscodestyle
   \let\\=\@normalcr
   \let\hspre\(\let\hspost\)%
   \pboxed}%
  {\endpboxed\)%
   \par\noindent
   \ignorespacesafterend}

\newcommand{\compaths}{\sethscode{compathscode}}

% "plain" mode is the proposed default.
% It should now work with \centering.
% This required some changes. The old version
% is still available for reference as oldplainhscode.

\newenvironment{plainhscode}%
  {\hsnewpar\abovedisplayskip
   \advance\leftskip\mathindent
   \hscodestyle
   \let\hspre\(\let\hspost\)%
   \pboxed}%
  {\endpboxed%
   \hsnewpar\belowdisplayskip
   \ignorespacesafterend}

\newenvironment{oldplainhscode}%
  {\hsnewpar\abovedisplayskip
   \advance\leftskip\mathindent
   \hscodestyle
   \let\\=\@normalcr
   \(\pboxed}%
  {\endpboxed\)%
   \hsnewpar\belowdisplayskip
   \ignorespacesafterend}

% Here, we make plainhscode the default environment.

\newcommand{\plainhs}{\sethscode{plainhscode}}
\newcommand{\oldplainhs}{\sethscode{oldplainhscode}}
\plainhs

% The arrayhscode is like plain, but makes use of polytable's
% parray environment which disallows page breaks in code blocks.

\newenvironment{arrayhscode}%
  {\hsnewpar\abovedisplayskip
   \advance\leftskip\mathindent
   \hscodestyle
   \let\\=\@normalcr
   \(\parray}%
  {\endparray\)%
   \hsnewpar\belowdisplayskip
   \ignorespacesafterend}

\newcommand{\arrayhs}{\sethscode{arrayhscode}}

% The mathhscode environment also makes use of polytable's parray 
% environment. It is supposed to be used only inside math mode 
% (I used it to typeset the type rules in my thesis).

\newenvironment{mathhscode}%
  {\parray}{\endparray}

\newcommand{\mathhs}{\sethscode{mathhscode}}

% texths is similar to mathhs, but works in text mode.

\newenvironment{texthscode}%
  {\(\parray}{\endparray\)}

\newcommand{\texths}{\sethscode{texthscode}}

% The framed environment places code in a framed box.

\def\codeframewidth{\arrayrulewidth}
\RequirePackage{calc}

\newenvironment{framedhscode}%
  {\parskip=\abovedisplayskip\par\noindent
   \hscodestyle
   \arrayrulewidth=\codeframewidth
   \tabular{@{}|p{\linewidth-2\arraycolsep-2\arrayrulewidth-2pt}|@{}}%
   \hline\framedhslinecorrect\\{-1.5ex}%
   \let\endoflinesave=\\
   \let\\=\@normalcr
   \(\pboxed}%
  {\endpboxed\)%
   \framedhslinecorrect\endoflinesave{.5ex}\hline
   \endtabular
   \parskip=\belowdisplayskip\par\noindent
   \ignorespacesafterend}

\newcommand{\framedhslinecorrect}[2]%
  {#1[#2]}

\newcommand{\framedhs}{\sethscode{framedhscode}}

% The inlinehscode environment is an experimental environment
% that can be used to typeset displayed code inline.

\newenvironment{inlinehscode}%
  {\(\def\column##1##2{}%
   \let\>\undefined\let\<\undefined\let\\\undefined
   \newcommand\>[1][]{}\newcommand\<[1][]{}\newcommand\\[1][]{}%
   \def\fromto##1##2##3{##3}%
   \def\nextline{}}{\) }%

\newcommand{\inlinehs}{\sethscode{inlinehscode}}

% The joincode environment is a separate environment that
% can be used to surround and thereby connect multiple code
% blocks.

\newenvironment{joincode}%
  {\let\orighscode=\hscode
   \let\origendhscode=\endhscode
   \def\endhscode{\def\hscode{\endgroup\def\@currenvir{hscode}\\}\begingroup}
   %\let\SaveRestoreHook=\empty
   %\let\ColumnHook=\empty
   %\let\resethooks=\empty
   \orighscode\def\hscode{\endgroup\def\@currenvir{hscode}}}%
  {\origendhscode
   \global\let\hscode=\orighscode
   \global\let\endhscode=\origendhscode}%

\makeatother
\EndFmtInput
%
%% ----------------------
%% Stuff for TimCore


%% 'g' inferences




%% --------- Effects ------------


%% ----------------------




%% ----------- Applications --------------
%% %format applyx a b = a "@" b
%% %format applyx a b = a "\," b

%% Function application with no arguments

%% Function application with 1 argument

%% Function application with 1 argument, without parenthesis

%% Function application with many arguments

%% ----------- End of Applications --------------


%% Sequences e1; ...; en; v

%% xs etc stole syntax used in examples.

%% Old version: | for BNF, [] for choice
%%%format choice = "[ \! ]"
%%%format `choice` = "\mathop{[ \! ]}"
%% 1mm is about 0.3em with the current font

%% New version: tall | for BNF, fat | for choice
%%format vbar = "\mathop{\bigm|}"
%%format choice = "\boldsymbol{|}"
%%format `choice` = "\mathop{\boldsymbol{|}}"




%% Arrays and tuples
%% %format arrKW = "\textbf{arr}"
%% %format arr (x) = arrKW "\{" x "\}"



%% %format defKW = "\textbf{def}"

% only used in comments, but still prettier this way:
%%%%format map = "\textbf{map}"

%% Effects checking, like succeeds{e}


%% Unification
%%%format == = "\!=\!"


% format := = "\mapsto"
%% %format squiggt = "\!\leadsto\!"
%  format effs x = "\!\!<\!\!" x "\!\!>\!\!" dots


































%% ----------- Metatheory --------------

\section{Updateable references} \label{sec:store}

\begin{figure}
  $$
  \begin{array}{l}
    \textbf{Syntax extension} \\
    \begin{array}{lrcl}
      \text{References}         & r \\
      \text{Expressions}        & e   & ::= & \cdots \vb \ensuremath{\textbf{store}\;\Varid{h}\;\textbf{in}\,\{\Varid{e}\}} \\
      \text{Primops}            & op  & ::= & \cdots \vb \ensuremath{\textbf{alloc}} \vb \ensuremath{\textbf{read}} \vb \ensuremath{\textbf{write}} \\
      \text{Head values}        & hnf & ::= & \cdots \vb r \\
      \text{Execution contexts} & X   & ::= & \cdots \vb \ensuremath{\textbf{store}\;\Varid{h}\;\textbf{in}\,\{X\}} \\
      \text{Scope contexts}     & SC  & ::= & \cdots \vb \ensuremath{\textbf{store}\;\Varid{h}\;\textbf{in}\,\{\Conid{SC}\}} \\
      \text{Heap}               & h   & ::= & \ensuremath{\epsilon} \vb \ensuremath{\Varid{r}\mapsto\!\Varid{v},\Varid{h}} \\
      \text{Heap context}       & H   & ::= & \hole\ensuremath{,\Varid{h}} \vb \ensuremath{\Varid{r}\mapsto\!\Varid{v},\Conid{H}} \\
      \text{Store contexts}     & S   & ::= & \hole \vb \ensuremath{\Varid{v}\hspace{-0.14em}=\hspace{-0.14em}\Conid{S}} \vb \ensuremath{\Conid{S};\,\Varid{e}} \vb \ensuremath{se;\,\Conid{S}} \vb \ensuremath{\exists\Varid{x}.\,\Conid{S}} \\
      \text{Store-op free exprs}&\ensuremath{se} & ::= & \ensuremath{\Varid{v}} \vb \ensuremath{se_1\mathrel{=}se_2} \vb \ensuremath{se_1;\,se_2} \vb \ensuremath{\exists\Varid{x}.\,se}
                                              \vb \ensuremath{\Varid{sp}(\Varid{v})}
                                              %%% need sv \vb |split sf sv_1 sv_2|
                                              \\
      \text{Results}            & w   & ::= & \ensuremath{\Varid{v}} \vb \ensuremath{\Varid{v}\hspace{0.2em}\mathop{\rule[-0.1em]{0.15em}{0.75em}}\hspace{0.2em}\Varid{e}} \\
      \text{Non-store primops}  & sp  & ::= & \text{any, except } \ensuremath{\textbf{alloc}}, \ensuremath{\textbf{read}}, \ensuremath{\textbf{write}} \\
      \text{Non-store expression} & oe  & ::= & \text{like \ensuremath{\Varid{e}}, but not \ensuremath{\Varid{w}}, \ensuremath{\textbf{store}}, or \ensuremath{\textbf{fail}}} \\
    \end{array} \\
\\

  \textbf{Axiom extensions} \\
  \textit{Normalization change} \\
  \begin{array}{l r @{\;}c@{\;} l l}
%  \begin{array}{lrcll}
    \rulename{exi-float} & \ensuremath{X [\mskip1.5mu \exists\Varid{x}.\,\Varid{e}\mskip1.5mu]} & \movesto & \ensuremath{\exists\Varid{x}.\,X [\mskip1.5mu \Varid{e}\mskip1.5mu]} & \text{if $\ensuremath{X} \ne \hole,\,  x \notin \freevars{\ensuremath{X}}$, \alphamark}\\
    &&&& \text{if there is \ensuremath{\textbf{store}} in \ensuremath{X} then \ensuremath{\Varid{e}} $\in$ \ensuremath{\Varid{ce}}} \\
  \end{array} \\

  \textit{Reference ops} \\
  \begin{array}{l r @{\;}c@{\;} l l}
%  \begin{array}{lrcll}
    \rulename{ref-alloc} & \ensuremath{\textbf{store}\;\Varid{h}\;\textbf{in}\,\{\Conid{S} [\mskip1.5mu \textbf{alloc}(\Varid{v})\mskip1.5mu]\}} & \movesto & \ensuremath{\textbf{store}\;\Varid{r}\mapsto\!\Varid{v},\Varid{h}\;\textbf{in}\,\{\Conid{S} [\mskip1.5mu \Varid{r}\mskip1.5mu]\}} \\
    &&& \quad \freevars{\ensuremath{\Varid{v}}} \# \boundvars{\ensuremath{\Conid{S}}}, \text{\ensuremath{\Varid{r}} fresh}  \\
    \rulename{ref-read} & \ensuremath{\textbf{store}\;\Conid{H} [\mskip1.5mu \Varid{r}\mapsto\!\Varid{v}\mskip1.5mu]\;\textbf{in}\,\{\Conid{S} [\mskip1.5mu \textbf{read}(\Varid{r})\mskip1.5mu]\}} & \movesto & \ensuremath{\textbf{store}\;\Conid{H} [\mskip1.5mu \Varid{r}\mapsto\!\Varid{v}\mskip1.5mu]\;\textbf{in}\,\{\Conid{S} [\mskip1.5mu \Varid{v}\mskip1.5mu]\}} \\
    &&& \quad \freevars{\ensuremath{\Varid{v}}} \# \boundvars{\ensuremath{\Conid{S}}}, \text{\alphamark}  \\
    \rulename{ref-write} &
        \ensuremath{\textbf{store}\;\Conid{H} [\mskip1.5mu \Varid{r}\mapsto\!\Varid{v}_{\mathrm{1}}\mskip1.5mu]\;\textbf{in}\,\{\Conid{S} [\mskip1.5mu \textbf{write}\langle\Varid{r},\Varid{v}_{\mathrm{2}}\rangle\mskip1.5mu]\}}
        & \movesto & \ensuremath{\textbf{store}\;\Conid{H} [\mskip1.5mu \Varid{r}\mapsto\!\Varid{v}_{\mathrm{2}}\mskip1.5mu]\;\textbf{in}\,\{\Conid{S} [\mskip1.5mu \langle\rangle\mskip1.5mu]\}} \\
    &&& \quad \freevars{\ensuremath{\Varid{v}_{\mathrm{2}}}} \# \boundvars{\ensuremath{\Conid{S}}}  \\
  \end{array} \\
    \\[-0.6em]

  \textit{Store duplication} \\
  \begin{array}{l r @{\;}c@{\;} l l}
%  \begin{array}{lrcll}
    \rulename{st-split-dup} & \hspace{-1em} \ensuremath{\textbf{store}\;\Varid{h}\;\textbf{in}\,\{\Conid{S} [\mskip1.5mu \textbf{split}(\Varid{oe})\{\Varid{f},\Varid{g}\}\mskip1.5mu]\}} & \movesto & \ensuremath{\textbf{store}\;\Varid{h}\;\textbf{in}\,\{\Conid{S} [\mskip1.5mu \textbf{split}(\textbf{store}\;\Varid{h}\;\textbf{in}\,\{\Varid{oe}\})\{\Varid{f},\Varid{g}\}\mskip1.5mu]\}} \\
    &&& \quad \freevars{\ensuremath{\Varid{h}}} \# \boundvars{\ensuremath{\Conid{S}}}, \text{\alphamark}  \\
    \rulename{st-choice-dup} & \ensuremath{\textbf{store}\;\Varid{h}\;\textbf{in}\,\{\Varid{oe}\hspace{0.2em}\mathop{\rule[-0.1em]{0.15em}{0.75em}}\hspace{0.2em}\Varid{e}\}} & \movesto & \ensuremath{\textbf{store}\;\Varid{h}\;\textbf{in}\,\{\textbf{store}\;\Varid{h}\;\textbf{in}\,\{\Varid{oe}\}\hspace{0.2em}\mathop{\rule[-0.1em]{0.15em}{0.75em}}\hspace{0.2em}\Varid{e}\}} \\
  \end{array} \\

    \\[-0.6em]
  \textit{Store commit} \\
  \begin{array}{l r @{\;}c@{\;} l l}
%  \begin{array}{lrcll}
     \rulename{st-split} & \ensuremath{\textbf{store}\;\Varid{h}_{\mathrm{1}}\;\textbf{in}\,\{\Conid{S} [\mskip1.5mu \textbf{split}(\textbf{store}\;\Varid{h}_{\mathrm{2}}\;\textbf{in}\,\{\Varid{w}\})\{\Varid{f},\Varid{g}\}\mskip1.5mu]\}} & \movesto &
         \ensuremath{\textbf{store}\;\Varid{h}_{\mathrm{2}}\;\textbf{in}\,\{\Conid{S} [\mskip1.5mu \textbf{split}(\Varid{w})\{\Varid{f},\Varid{g}\}\mskip1.5mu]\}} \\
    &&& \quad \freevars{\ensuremath{\Varid{h}_{\mathrm{2}}}} \# \boundvars{\ensuremath{\Conid{S}}}  \\
     \rulename{st-choice} & \ensuremath{\textbf{store}\;\Varid{h}_{\mathrm{1}}\;\textbf{in}\,\{\Conid{S} [\mskip1.5mu (\textbf{store}\;\Varid{h}_{\mathrm{2}}\;\textbf{in}\,\{\Varid{w}\})\hspace{0.2em}\mathop{\rule[-0.1em]{0.15em}{0.75em}}\hspace{0.2em}\Varid{e}\mskip1.5mu]\}} & \movesto & \ensuremath{\textbf{store}\;\Varid{h}_{\mathrm{2}}\;\textbf{in}\,\{\Conid{S} [\mskip1.5mu \Varid{w}\hspace{0.2em}\mathop{\rule[-0.1em]{0.15em}{0.75em}}\hspace{0.2em}\Varid{e}\mskip1.5mu]\}} \\
    &&& \quad \freevars{\ensuremath{\Varid{h}_{\mathrm{2}}}} \# \boundvars{\ensuremath{\Conid{S}}}  \\
  \end{array} \\

    \\[-0.6em]
    \textit{Unification} \\
    \begin{array}{lrcll}
      \text{Extension with the obvious axioms making equal references unify,
        and anything else fail.}
    \end{array} \\
  \textit{Top level} \\
  \begin{array}{lrcll}
    \text{Start top level reduction of \ensuremath{\Varid{e}} with \ensuremath{\textbf{store}\;\epsilon\;\textbf{in}\,\{\Varid{e}\}}.}
  \end{array} \\

  \end{array}
   $$
  \caption{The Verse calculus: store axioms} \label{fig:store-axioms}
  \Description{The Verse calculus: store axioms}
\end{figure}
The full Verse language has updatable references (à la ML).
There are three new primitive operations, \ensuremath{\textbf{alloc}}, \ensuremath{\textbf{read}}, and \ensuremath{\textbf{write}}.
The \ensuremath{\textbf{alloc}} creates a new reference with an initial value,
\ensuremath{\textbf{read}} extracts the value from a reference,
and \ensuremath{\textbf{write}} sets the value of a referene.

Modifying these references is transactional in the sense that if a computation fails,
then any updates will not be visible outside the construct that handles the failures.
\Eg,\\
\ensuremath{\Varid{r}\mathrel{\coloneqq}\textbf{alloc}(\mathrm{0});\,(\textbf{if}\;(\textbf{write}\langle\Varid{r},\mathrm{1}\rangle;\,\textbf{fail})\;\textbf{then}\;\mathrm{1}\;\textbf{else}\;\mathrm{2});\,\textbf{read}(\Varid{r})}\\
will have the value \ensuremath{\mathrm{0}},
because the \ensuremath{\textbf{write}} is part of an expression that fails, and so its effect is not visible.


To add updateable references we extend the system with syntax and rules from figure~\ref{fig:store-axioms}.
The \ensuremath{\textbf{store}\;\Varid{h}\;\textbf{in}\,\{\Varid{e}\}} indicates that \ensuremath{\Varid{e}} should be reduced using the heap \ensuremath{\Varid{h}}.
A heap is simply a mapping from references to values (one mapping being \ensuremath{\Varid{r}\mapsto\!\Varid{v}}).
A reference is some opaque type that supports equality (unification) and creation of a new reference.

The interaction of the new primitives with the store can be seen from the axioms.
The \ensuremath{\textbf{alloc}(\Varid{v})} operation creates a new reference and adds a binding with \ensuremath{\Varid{v}} to the store.
The \ensuremath{\textbf{read}(\Varid{r})} operation retrieves the value for reference \ensuremath{\Varid{r}} from the store,
and \ensuremath{\textbf{write}\langle\Varid{r},\Varid{v}\rangle} updates the reference \ensuremath{\Varid{r}} with \ensuremath{\Varid{v}} in the store.
All of these operations use the context \ensuremath{\Conid{S}} which ensures that there are no store operations
to the left of the hole, \ie, a store operation in the hole is the next one that should execute.

The interesting rules involve choice and \ensuremath{\textbf{split}} because store operations are transactional in the sense
that when an expressions fails, none of its store operations will happen.

When reducing \ensuremath{\textbf{split}(\Varid{e})\{\Varid{f},\Varid{g}\}} in an \ensuremath{\Conid{S}} hole, rule \rulename{st-split-dup}, the store is duplicated.
Any store operations inside the \ensuremath{\textbf{split}} will happen in this local copy of the store.
Note the two occurrences of \ensuremath{\Varid{h}} in the right hand side of \rulename{st-split-dup}.
If the reduction of \ensuremath{\Varid{e}} results in \ensuremath{\textbf{fail}} then rule \rulename{fail-elim} is used,
and the store from the failing computation is simply thrown away.
If the reduction of \ensuremath{\Varid{e}} results in a value (with or without more alternatives) then rule
\rulename{st-split} is used.  This rule replaces the outer store with the inner store,
since we know the inner computation has succeeded.

Similarely, the reduction of \ensuremath{\Varid{e}_{\mathrm{1}}\hspace{0.2em}\mathop{\rule[-0.1em]{0.15em}{0.75em}}\hspace{0.2em}\Varid{e}_{\mathrm{2}}} will duplicate the store into the first branch, \rulename{st-choice-dup}.
Here \ensuremath{\Varid{e}_{\mathrm{1}}} must not contain any store operation nor be a value.
And again, similarely, \rulename{st-choice} commits the new store and throws away the old.

The use of \ensuremath{\Varid{oe}} in the rules is to ensure that the rules cannot get stuck in a loop.
Using \ensuremath{\Varid{e}} instead of \ensuremath{\Varid{oe}} would mean that failing or committing would make the expression match the duplication rule again.
It also prevents the duplication rule from repeatedly duplicating the \ensuremath{\textbf{store}}.

Note that \ensuremath{\textbf{store}} is part of the \ensuremath{X} context, which means that the \ensuremath{\textbf{store}} can float inside
existentials.
This is necessary for the store rules to fire since the \ensuremath{\Conid{S}} context does not allow
going under existentials.
%If |S| contained |storeKW| then an |alloc| or |write| could update the store with
%a variable that is not in scope.

The semantics of \ensuremath{\textbf{for}(\Varid{d})\,\textbf{do}\,\Varid{e}} with respect to store effects is somewhat intricate.
The expression \ensuremath{\Varid{d}} is possibly multi-valued; any effects that happens when computing the
first value of \ensuremath{\Varid{d}} will be visible the first time \ensuremath{\Varid{e}} is computed.
Both these effects are then visible when computing the second value of \ensuremath{\Varid{d}}, and so on.
If any iteration of \ensuremath{\Varid{d}} fails, then the effects of that computation are not visible outside \ensuremath{\Varid{d}}.
This means that the desugaring of \ensuremath{\textbf{for}} into \ensuremath{\textbf{split}} needs to be more elaborate.

\ensuremath{\textbf{for}(\exists\Varid{x}_{\mathrm{1}} \mydots \Varid{x}_{\Varid{n}}.\,\Varid{d})\,\textbf{do}\,\Varid{e}}\\
means\\
\hspace*{5mm}\ensuremath{\Varid{f}\,\langle\rangle\mathrel{\coloneqq}\langle\rangle;\,}\\
\hspace*{5mm}\ensuremath{\Varid{g} (\Varid{v}) (\Varid{r})\mathrel{\coloneqq}(\Varid{v}\hspace{-0.14em}=\hspace{-0.14em}\langle{}\Varid{x}_{\mathrm{1}},\mydots,\Varid{x}_{\Varid{n}}\rangle;\,\Varid{cons}\langle\Varid{e},\textbf{split}(\Varid{r}\,\langle\rangle)\{\Varid{f},\Varid{g}\}\rangle;\,}\\
\hspace*{5mm}\ensuremath{\textbf{split}(\exists\Varid{x}_{\mathrm{1}}\;\mydots\;\Varid{x}_{\Varid{n}}.\,\Varid{d};\,\langle{}\Varid{x}_{\mathrm{1}},\mydots,\Varid{x}_{\Varid{n}}\rangle)\{\Varid{f},\Varid{g}\}}

To support limited store operations ({\eg}, \ensuremath{\textbf{read}},
but not \ensuremath{\textbf{write}}) we can equip the store with a set of currently allowed operations.
We also need some extra primitives that modify this set.

\subsection{Examples}
\lennart{Not yet}

